\documentclass{article}

\usepackage[utf8]{inputenc}
\usepackage[T1]{fontenc}
\usepackage[spanish]{babel}
\usepackage{amsmath}
\usepackage{amssymb}
\usepackage{array}
\usepackage{geometry}
\usepackage{enumitem}

\geometry{margin=2.5cm}

\begin{document}

\section{Guía de comprensión de Flex}

Esta guía tiene como objetivo comprobar la comprensión del funcionamiento
básico de Flex. Las preguntas deben responderse con palabras propias,
utilizando ejemplos cuando sea necesario.

\subsection{Función de Flex}

\begin{enumerate}
    \item ¿Qué problema resuelve Flex dentro de un compilador?

    \item ¿Qué recibe como entrada Flex?

    \item ¿Qué produce como resultado?

    \item ¿Cuál es la diferencia entre lo que hace Flex y lo que hará
    posteriormente Bison?
\end{enumerate}

\subsection{Lexemas y tokens}

\begin{enumerate}
    \item ¿Qué es un lexema?

    \item ¿Qué es un token?

    \item Explicar con un ejemplo la diferencia entre ambos.

    Por ejemplo, para:

\begin{verbatim}
int edad = 123;
\end{verbatim}

    identificar los lexemas:

\begin{verbatim}
int
edad
=
123
;
\end{verbatim}

    y asignarles sus respectivos tokens.

    \item ¿Se puede decir que un token es una clasificación del lexema?
    Explicar por qué.
\end{enumerate}

\subsection{Patrones}

Considérense las siguientes reglas:

\begin{verbatim}
"int"                         { return INT; }
[0-9]+                        { return NUMERO; }
[a-zA-Z_][a-zA-Z0-9_]*        { return IDENTIFICADOR; }
"+"                           { return SUMA; }
\end{verbatim}

\begin{enumerate}
    \item ¿Qué es cada parte de una regla como:

\begin{verbatim}
[0-9]+ { return NUMERO; }
\end{verbatim}

    Especificar qué representa:

\begin{verbatim}
[0-9]+
return NUMERO
\end{verbatim}

    \item ¿Qué relación existe entre un patrón y un lexema?

    \item ¿Flex busca palabras concretas o patrones que puedan reconocer
    partes de la entrada?
\end{enumerate}

\subsection{Cómo encuentra un emparejamiento}

Considérense las reglas:

\begin{verbatim}
[0-9]     { return DIGITO; }
[0-9]+    { return NUMERO; }
\end{verbatim}

y la entrada:

\begin{verbatim}
1234 + 5
\end{verbatim}

\begin{enumerate}
    \item Cuando Flex está en el primer carácter, ¿qué intenta hacer?

    \item ¿Qué coincidencias puede encontrar?

    \item ¿Por qué no devuelve inmediatamente el primer carácter
    reconocido?

    \item ¿Qué significa que Flex busque el emparejamiento más largo?

    \item ¿Qué lexema termina reconociendo?

    \item ¿Qué ocurre con la posición de análisis después de reconocerlo?
\end{enumerate}

\subsection{Maximal munch}

Explicar con palabras propias qué significa el principio de
\textit{emparejamiento más largo} o \textit{maximal munch}.

Considérense las reglas:

\begin{verbatim}
"="      { return ASIGNACION; }
"=="     { return IGUALDAD; }
\end{verbatim}

y la entrada:

\begin{verbatim}
==
\end{verbatim}

\begin{enumerate}
    \item ¿Qué patrones coinciden?

    \item ¿Cuál gana?

    \item ¿Por qué?
\end{enumerate}

\subsection{Empates entre patrones}

Considérense las reglas:

\begin{verbatim}
"if"    { return IF; }
[a-z]+  { return IDENTIFICADOR; }
\end{verbatim}

y la entrada:

\begin{verbatim}
if
\end{verbatim}

\begin{enumerate}
    \item ¿Qué reglas coinciden?

    \item ¿Qué longitud tiene cada emparejamiento?

    \item ¿Cuál gana?

    \item ¿Qué ocurre cuando dos reglas producen un emparejamiento de
    la misma longitud?

    \item ¿Qué cambiaría si se invirtiera el orden de las reglas?
\end{enumerate}

\subsection{\texttt{yytext}}

\begin{enumerate}
    \item ¿Qué es \texttt{yytext}?

    \item Si Flex reconoce:

\begin{verbatim}
1234
\end{verbatim}

    ¿qué contiene \texttt{yytext}?

    \item Si posteriormente reconoce:

\begin{verbatim}
+
\end{verbatim}

    ¿qué contiene \texttt{yytext}?

    \item ¿\texttt{yytext} almacena todos los lexemas encontrados
    anteriormente?

    \item ¿Para qué podría ser útil \texttt{yytext} en un compilador?
    Considerar especialmente números, identificadores y errores léxicos.
\end{enumerate}

\subsection{\texttt{yyleng}}

\begin{enumerate}
    \item ¿Qué contiene \texttt{yyleng}?

    \item Si:

\begin{verbatim}
yytext = "1234"
\end{verbatim}

    ¿cuánto vale \texttt{yyleng}?

    \item ¿Cuál es la diferencia entre:

\begin{verbatim}
yytext
yyleng
yylex()
\end{verbatim}
\end{enumerate}

\subsection{\texttt{yylex()}}

\begin{enumerate}
    \item ¿Qué hace \texttt{yylex()}?

    \item ¿Qué devuelve?

    \item ¿Cuál es la diferencia entre \texttt{yytext},
    \texttt{yyleng} y \texttt{yylex()}?

    \item Cuando usamos Flex junto con Bison, ¿quién llama normalmente
    a \texttt{yylex()}?
\end{enumerate}

\subsection{Entrada no reconocida}

Considérense las siguientes reglas:

\begin{verbatim}
"int"       { return INT; }
[0-9]+      { return NUMERO; }
"+"         { return SUMA; }

. {
    printf("Error léxico: %s\n", yytext);
}
\end{verbatim}

y la entrada:

\begin{verbatim}
int x = 10 @ 20
\end{verbatim}

\begin{enumerate}
    \item ¿Qué ocurre con \texttt{int}?

    \item ¿Qué ocurre con \texttt{10}?

    \item ¿Qué ocurre con \texttt{@}?

    \item ¿Qué contiene \texttt{yytext} cuando se ejecuta la regla
    \texttt{.}?

    \item ¿Qué ocurre después de procesar \texttt{@}?
\end{enumerate}

\subsection{Espacios}

Considérese la regla:

\begin{verbatim}
[ \t\n]+    { /* ignorar */ }
\end{verbatim}

\begin{enumerate}
    \item ¿Qué reconoce esta regla?

    \item ¿Por qué normalmente no queremos generar tokens para los
    espacios?

    \item ¿Qué ocurre con una entrada como:

\begin{verbatim}
int     x
\end{verbatim}

    ¿Flex procesa los espacios? ¿Los devuelve como tokens?
\end{enumerate}

\subsection{Ejercicio integrador}

Considérense las siguientes reglas:

\begin{verbatim}
"if"                          { return IF; }
"=="                          { return IGUALDAD; }
"="                           { return ASIGNACION; }
[0-9]+                        { return NUMERO; }
[a-zA-Z_][a-zA-Z0-9_]*        { return IDENTIFICADOR; }
"+"                           { return SUMA; }
[ \t\n]+                      { /* ignorar */ }
.                             { return ERROR; }
\end{verbatim}

Analizar la siguiente entrada:

\begin{verbatim}
if edad == 123 + edad @ 5
\end{verbatim}

Completar la siguiente tabla:

\[
\begin{array}{c|c|c|c}
\text{Entrada restante} & \texttt{yytext} &
\texttt{yyleng} & \text{Token} \\
\hline
\texttt{if edad == 123 + edad @ 5} & & & \\
\texttt{edad == 123 + edad @ 5} & & & \\
\texttt{== 123 + edad @ 5} & & & \\
\texttt{123 + edad @ 5} & & & \\
\texttt{+ edad @ 5} & & & \\
\texttt{edad @ 5} & & & \\
\texttt{@ 5} & & & \\
\texttt{5} & & &
\end{array}
\]

No es necesario incluir las posiciones correspondientes a los espacios
que son ignorados.

\subsection{Pregunta final}

Explicar con palabras propias:

\begin{quote}
¿Qué hace Flex desde que recibe la entrada
\texttt{1234 + 56} hasta que termina produciendo los tokens?
\end{quote}

La explicación debería incluir el reconocimiento de patrones, el
emparejamiento más largo, el uso de \texttt{yytext} y \texttt{yyleng},
la ejecución de las acciones y la continuación del análisis sobre la
entrada restante.

\end{document}